\section{Quantitative estimates}\label{sec:estimates}

This appendix proves the estimates used in \cref{sec:comparison,sec:support}. The required comparisons come from Chi, Sun, and Wang~\cite[Sections 3--4 and Appendix A]{CSW2026}; the algebraic verifications below use explicit polynomial inequalities and exact rational certificates. Throughout, $c=c_k$ is determined by $g_k(c)=1$, and $x=1/(k+c)$, $a=kx$, and $\delta=1-a=cx$. In particular, at the critical mean we have
\begin{align*}
p_0\coloneqq(1-x)^k=1-a^k.
\end{align*}
The functions $A$ and $H$, their nodes $W_t$, and the interpolated values $e_j^*$ and $f_j^*$ are those of \cref{sec:comparison}.

\subsection{Bounds on the critical mean}

We first locate $c_k$ in a short interval. The lower bound will be used both in the comparison of the first four nodes and in the estimates for the conditional variance.

\begin{lemma}\label{lem:crossing-bounds}
For every integer $k\ge2$,
\begin{align}\label{eq:appendix-crossing}
c_k>\frac{11}{24}+\frac{2}{25k}.
\end{align}
Moreover, $c_2=1/2$ and $c_k<1/2$ for $k\ge3$. If $k\ge6$, then
\begin{align}\label{eq:appendix-elementary-bounds}
\frac{12}{13}<a<1,\qquad ac>\frac5{12},\qquad \frac13<p_0<\mathrm e^{-1}<\frac38.
\end{align}
\end{lemma}

\begin{proof}
The assertions $c_2=1/2$ and $c_k<1/2$ for $k\ge3$ follow from the calculation of $g_k(1/2)$ in \cref{sec:reduction}. Since $g_k$ is strictly increasing, to prove \cref{eq:appendix-crossing} it suffices to show $g_k(c_*)<1$, where
\begin{align*}
c_*\coloneqq\frac{11}{24}+\frac{2}{25k},\qquad
\eta\coloneqq\frac{1/2-c_*}{k}=\frac1{24k}-\frac2{25k^2},\qquad b\coloneqq1-\eta.
\end{align*}
Here $0<\eta<1/(24k)$. Expanding the difference of the two powers about $b$ gives
\begin{align*}
g_k(c_*)&=\left(b+\frac1{2k}\right)^k-\left(b-\frac1{2k}\right)^k\\
&=b^{k-1}\sum_{j\ge0}\frac{\binom{k-1}{2j}}{2j+1}(2kb)^{-2j}.
\end{align*}
The sum is finite, with the convention that a binomial coefficient is zero outside its usual range. For every $j\ge1$, pair the factors in the numerator of $\binom{k-1}{2j}$ and use $(2h)(2h+1)\ge6h$ for $h\in[j]$. This gives
\begin{align*}
\frac{\binom{k-1}{2j}}{2j+1}
\le \frac{((k-1)(k-2))^j}{6^j j!}.
\end{align*}
Consequently,
\begin{align*}
g_k(c_*)\le b^{k-1}\exp\left(\frac{(k-1)(k-2)}{24k^2b^2}\right).
\end{align*}
For $0<\eta<1$, the inequality $(1-\eta)^2\log(1-\eta)<-\eta+3\eta^2/2$ follows by differentiation from $-\log(1-\eta)<\eta/(1-\eta)$. Applying it to $b=1-\eta$, we obtain
\begin{align*}
\frac{b^2\log g_k(c_*)}{k-1}
&\le b^2\log b+\frac{k-2}{24k^2}\\
&< -\eta+\frac32\eta^2+\frac{k-2}{24k^2}\\
&<\frac1{k^2}\left(\frac2{25}-\frac1{12}+\frac1{384}\right)<0.
\end{align*}
Thus $g_k(c_*)<1$, as required.

Now suppose $k\ge6$. Since $c<1/2$, we have $a>k/(k+1/2)\ge12/13$. Combining this with $c>11/24$ gives $ac>11/26>5/12$. Also, $x<1/k$ and the sequence $(1-1/k)^k$ is increasing, so
\begin{align*}
p_0>\left(1-\frac1k\right)^k\ge\left(\frac56\right)^6>\frac13.
\end{align*}
For the other direction, $x>1/(k+1/2)$ implies
\begin{align*}
p_0<\left(\frac{2k-1}{2k+1}\right)^k<\mathrm e^{-1}.
\end{align*}
The last inequality follows from $\log((1+u)/(1-u))>2u$ with $u=1/(2k)$. Finally, the exponential series gives $\mathrm e>1+1+1/2+1/6=8/3$, proving $\mathrm e^{-1}<3/8$.
\end{proof}

\subsection{The first four comparison nodes}

We next verify the node inequalities needed to restrict the possible value of $w$. The strict inequality at the fourth node will also be used later when we compare $r^*$ and $M^*$.

\begin{lemma}\label{lem:node-estimates}
At $x=\beta_k$, the comparison functions satisfy
\begin{align}\label{eq:appendix-endpoint-nodes}
A(W_1)=H(W_1)=0,\qquad A(W_k)=H(W_k),
\end{align}
and
\begin{align}\label{eq:appendix-strict-nodes}
A(W_t)<H(W_t)\qquad\text{for }2\le t\le\min\{4,k-1\}.
\end{align}
If $k\ge5$, then
\begin{align}\label{eq:appendix-fourth-node-gap}
H(W_4)-A(W_4)>
\frac{c^4\prod_{h\in[4]}(1-h/k)}{96(1-c/5)}.
\end{align}
\end{lemma}

\begin{proof}
Substituting the node values $f_{j,t}$ into the definition of $A$ and setting $h=k-j$ gives
\begin{align}\label{eq:appendix-node-formula}
A(W_t)=\sum_{1\le h<t}\binom{k}{h}\left(\frac ck\right)^h
\left[1-\left(1-\frac1t\right)^{t-h}\right].
\end{align}
The first equality in \cref{eq:appendix-endpoint-nodes} is immediate. For $t=k$, include the term $h=0$ and apply the binomial theorem. The resulting sum is $g_k(c)=1$, so
\begin{align*}
A(W_k)=g_k(c)-1+\left(1-\frac1k\right)^k=H(W_k).
\end{align*}

For $t\ge2$, write $Q_t\coloneqq A(W_t)/H(W_t)$. Expanding \cref{eq:appendix-node-formula} at $t=2,3,4$ yields
\begin{align*}
Q_2&=2c,\\
Q_3&=\frac{15c}{8}+\frac9{16}\left(1-\frac1k\right)c^2,\\
Q_4&=\frac{148c}{81}+\frac{56}{81}\left(1-\frac1k\right)c^2
+\frac{32}{243}\left(1-\frac1k\right)\left(1-\frac2k\right)c^3.
\end{align*}
In their respective ranges $k\ge3$ and $k\ge4$, the differences are
\begin{align*}
Q_3-Q_2&=\frac c{16}\left[9\left(1-\frac1k\right)c-2\right]>0,\\
Q_4-Q_3&=\frac c{3888}\left[501\left(1-\frac1k\right)c
+512\left(1-\frac1k\right)\left(1-\frac2k\right)c^2-186\right]>0.
\end{align*}
To see the strict signs, it is enough to use $c>5/12$. For the first bracket, its value is then greater than $1/2$. For the second, substituting $k=4$ and $c=5/12$ gives $187/48>0$, and every positive term increases with $k$ and $c$. Since $Q_k=1$, these comparisons prove \cref{eq:appendix-strict-nodes} when $k\le4$.

It remains to treat $k\ge5$. Put $h=1/k$ and define
\begin{align*}
G(c,h)\coloneqq{}&1+\frac{1-h}{2}(2c-1)
+\frac{(1-h)(1-2h)}6(3c^2-3c+1)\\
&+\frac{(1-h)(1-2h)(1-3h)}{24}(4c^3-6c^2+4c-1),\\
A_4\coloneqq{}&\frac{37c}{64}+\frac7{32}(1-h)c^2
+\frac1{24}(1-h)(1-2h)c^3,\\
T_4\coloneqq{}&\frac{c^4(1-h)(1-2h)(1-3h)(1-4h)}{24(1-c/5)},\\
F_4\coloneqq{}&\frac{81}{256}-A_4-\frac{T_4}{4}.
\end{align*}
Here $A_4=A(W_4)$ and $81/256=H(W_4)$. Taylor's formula through degree four, applied to $(1+c/k)^k$ and $(1-(1-c)/k)^k$, gives $G(c,1/k)\le g_k(c)$. Indeed, the first fifth-order remainder is nonnegative and the second is nonpositive; the fourth-order difference is precisely $G(c,1/k)$.

To prove $F_4>0$, write $y=1/2-c$ and use the following polynomial identity:
\begin{align}\label{eq:appendix-fourth-identity}
F_4+\frac{4-h}{5(1-h)}(G(c,h)-1)
=\frac{1-4h}{3840(5-c)}
\left(R_0(c)+hR_1(c)+h^2R_2(c)+1200c^4h^3\right),
\end{align}
where
\begin{align*}
R_0(c)&=1+13y+y^2(1732-728y-552y^2),\\
R_1(c)&=48+y(114+1636y+1128y^2+2000y^3),\\
R_2(c)&=-240+c\left(1008-1632c+c^2(1248-2392c)\right).
\end{align*}
For $0\le y\le1/2$, the quadratic expression in $R_0$ is at least $1230$. Thus $R_0\ge1$ and $R_1\ge48$. For $0\le c\le1/2$, we have $1008-1632c\ge192$ and $1248-2392c\ge52$, so $R_2\ge-240$. Since $0<h\le1/5$, the bracket in \cref{eq:appendix-fourth-identity} is at least
\begin{align*}
1+48h-240h^2\ge1.
\end{align*}
The right-hand side of \cref{eq:appendix-fourth-identity} is therefore positive. At $c=c_k$, we have $G(c,1/k)\le g_k(c)=1$, and hence $F_4>0$. This is exactly \cref{eq:appendix-fourth-node-gap}. The inequalities $Q_2<Q_3<Q_4<1$ now give \cref{eq:appendix-strict-nodes} for every $k\ge5$.
\end{proof}

\subsection{The interpolation coefficients}

We now begin the proof of \cref{lem:comparison-estimates}. Fix $k\ge6$ and a segment $[W_t,W_{t+1}]$ with $4\le t<k$. The quantities $E^*,V^*,L,r_b,r^*,M^*$ are defined in \cref{eq:comparison-definitions}. In particular, $L=t(t-1)/(2k)$.

We first verify the nonnegativity of the coefficients in \cref{lem:interpolation} for the weights $\kappa_j=Lv_j$ and costs $c_j=a^kb_j/k$. Let $\alpha_{k-\ell}$ be the slope between the nodes $t,t+1$ of the curve $C_{k-\ell}$. Dividing the difference of the tail coordinates by the difference of the excess coordinates gives
\begin{align}\label{eq:appendix-slope}
\alpha_{k-\ell}
=\frac1{t(t+1)\left[1-(1-t^{-2})^{t-\ell}\right]}-\frac1t
\ge\frac{(2t+1)\ell-t-1}{2t(t+1)(t-\ell)}
\qquad(1\le\ell<t).
\end{align}
For completeness, the inequality follows from
\begin{align*}
\frac1{1-q^m}\ge\frac1{m(1-q)}+\frac{m-1}{2m}
\qquad(0<q<1,\ m\ge1).
\end{align*}
To prove this, sum $q^i+q^{m-i}\le1+q^m$ over $1\le i<m$ and rearrange the resulting bound on $(1-q^m)/(1-q)$. The inequality is strict for $m\ge2$. Taking $q=1-t^{-2}$ and $m=t-\ell$ proves \cref{eq:appendix-slope}, strictly when $\ell\le t-2$.

Let $\Gamma_j$ be the coefficients from \cref{lem:interpolation}, and set $h_\ell\coloneqq k\Gamma_{k-\ell}/v_{k-\ell}$. The initial coefficient gives $h_t=kL$. Substituting the ratios of $v_j$ and the costs $a^kb_j/k$ into the recurrence for $\Gamma_j$ gives
\begin{align}\label{eq:appendix-multiplier-recurrence}
(1+\alpha_{k-\ell})h_\ell
=\frac c\ell\left[\left(1-\frac\ell k\right)h_{\ell+1}-a\right]
+kL\alpha_{k-\ell}
\qquad(1\le\ell<t).
\end{align}
The lower bound in \cref{eq:appendix-slope} implies
\begin{align*}
(kL-1)\alpha_{k-\ell}&>\frac54 &&(3\le\ell<t),\\
(kL-2/3)\alpha_{k-2}&>\frac34,\\
kL\alpha_{k-1}&>\frac{t}{4(t+1)}\ge\frac15.
\end{align*}
For the first inequality, the right-hand side of \cref{eq:appendix-slope} increases with $\ell$, so it is enough to use $\ell=3$. The resulting bound is
\begin{align*}
\frac{(t-2)(5t+2)}{4t(t-3)}>\frac54.
\end{align*}
For the second inequality, the difference from $3/4$ of the resulting lower bound has numerator $3t^2+3t-4>0$ and positive denominator $12t(t+1)(t-2)$. The third inequality follows by substituting $\ell=1$ and using the strict part of \cref{eq:appendix-slope}.

Starting with $h_t=kL\ge6$, downward induction in \cref{eq:appendix-multiplier-recurrence} gives $h_\ell>1$ for $3\le\ell\le t$. Indeed, the difference of its right-hand side from $1+\alpha_{k-\ell}$ is greater than $-1/6+5/4-1>0$. The last two steps are
\begin{align*}
(1+\alpha_{k-2})(h_2-2/3)&>-\frac ck+\frac1{12}>0,\\
(1+\alpha_{k-1})h_1
&>\frac15+c\left[\frac23\left(1-\frac1k\right)-a\right]\\
&=\frac15-\frac c3+\frac ck\left(ac-\frac23\right)\\
&>\frac15-\frac16-\frac1{48}=\frac1{80}>0.
\end{align*}
Here we used $c<1/2$, $k\ge6$, and $ac>5/12$ from \cref{lem:crossing-bounds}. Thus all the required $\Gamma_j$ are positive. The weights themselves are increasing, because
\begin{align*}
\frac{v_{j+1}}{v_j}=\frac{k(k-1-j)}{c(j+1)}>1
\qquad(1\le j\le k-2).
\end{align*}

The estimates involving $r_b$ also follow directly from \cref{lem:crossing-bounds}. Since $H(W_t)\ge H(W_4)=81/256$, $p_0<3/8$, and $c>4/9$, we have
\begin{align*}
\frac{r_b}{x}-1
=\frac{(t+c)H(W_t)}{p_0}-1
>\frac{27t-20}{32}>0.
\end{align*}
In particular $r_b>x$. Moreover,
\begin{align}\label{eq:appendix-rb-ratio}
\frac{Lx}{(r_b-x)^2}
<\frac{1664t(t-1)}{3(27t-20)^2}
\le\frac{104}{121}<\frac9{10}.
\end{align}
For the middle inequality, the function $t(t-1)/(27t-20)^2$ decreases on $t\ge4$, since its derivative is $(20-13t)/(27t-20)^3$. Its value at $t=4$ gives the stated constant.

\subsection{Comparing the conditional deficit and the support bound}

It remains to prove $a^k>V^*$ and $r^*-(19/20)M^*>x/5$ in \cref{lem:comparison-estimates}. We first prove these inequalities at the nodes and then interpolate.

Fix a node index $t\ge4$, put $q=(t-1)/t$, and write $E_t\coloneqq E^*(W_t)$ and $V_t\coloneqq V^*(W_t)$. Let $N$ have the binomial distribution with parameters $k-1$ and $\delta$. The node formulas give
\begin{align}\label{eq:appendix-node-identities}
1-V_t&=\E q^{(t-1-N)_+},\\
k(\delta+E_t)&=(t-a)(1-V_t)+\E(N-t+1)_+.\nonumber
\end{align}
The first identity follows by viewing $N$ as the number of zero choices in the thinning formula. To check the second, take an independent Bernoulli variable $I$ of mean $\delta$ and set $J=N+I$. The formula for $e_{j,t}$ gives
\begin{align*}
kE_t=\E\left[tq^{(t-J)_+}-J+(J-t)_+\right].
\end{align*}
Conditional on $N$, the sum of the first and third terms has expectation
\begin{align*}
(t-a)q^{(t-1-N)_+}+(N-t+1)_+.
\end{align*}
Using $\E J=k\delta$ yields the second identity in \cref{eq:appendix-node-identities}.

Since $q^{(t-1-N)_+}\ge q^{t-1}$ and $(1-1/t)^{t-1}>\mathrm e^{-1}>p_0$, we have $1-V_t>p_0$ and hence $V_t<a^k$. Define
\begin{align*}
z_t\coloneqq\frac{a^k-V_t}{p_0}>0,\qquad
G_n(a,z)\coloneqq n+1-2a+(n+1-a)z-\frac{19a}{20z}.
\end{align*}
Substituting $1-V_t=p_0(1+z_t)$ into \cref{eq:appendix-node-identities}, and using $M^*(W_t)=x(1+1/z_t)$, gives
\begin{align}\label{eq:appendix-target-reduction}
k\left(r^*(W_t)-\frac{19}{20}M^*(W_t)-\frac{x}{20}\right)
=G_{t-1}(a,z_t)+\frac{\E(N-t+1)_+}{p_0}.
\end{align}
For fixed $a$, the function $G_n(a,z)$ increases with $z>0$. Thus it is enough to prove $G_{t-1}(a,z_t)>3a/20$.

At $t=4$, conditioning on $N$ as above also gives
\begin{align*}
(4-a)(1-V_4)-4p_0
=4a^k\bigl[H(W_4)-A(W_4)\bigr]-a\Pr(N\ge4).
\end{align*}
The binomial tail on the right is bounded by its first term and a geometric series. Indeed, after division by $a^k$ its consecutive terms have ratio at most $c/5$, so
\begin{align*}
\frac{a\Pr(N\ge4)}{a^k}
\le\frac{c^4\prod_{h\in[4]}(1-h/k)}{24(1-c/5)}
<4\bigl[H(W_4)-A(W_4)\bigr],
\end{align*}
where the last inequality is \cref{eq:appendix-fourth-node-gap}. Hence $(4-a)(1-V_4)>4p_0$, or equivalently $z_4>a/(4-a)$. It follows that
\begin{align*}
G_3(a,z_4)>G_3\left(a,\frac{a}{4-a}\right)=\frac{4-a}{20}>\frac{3a}{20}.
\end{align*}

Now suppose $5\le t<k$, and put $n=t-1\ge4$. Define the comparison value
\begin{align*}
\widehat z_t(u)\coloneqq\frac{u(19t+4u)}{20t(t-u)}.
\end{align*}
Direct substitution gives
\begin{align}\label{eq:appendix-zhat}
G_{t-1}(a,\widehat z_t(a))-\frac{3a}{20}
=\frac{a(t-4a)(t-a)}{5t(19t+4a)}>0.
\end{align}
We shall prove $z_t>\widehat z_t(a)$.

To estimate $1-V_t$, first omit the truncation in the exponent in \cref{eq:appendix-node-identities}. The binomial generating function then gives
\begin{align}\label{eq:appendix-untruncated-tail}
1-V_t&=H_n-T_n,\\
H_n&\coloneqq q^n(1+\delta/n)^{k-1},\nonumber\\
T_n&\coloneqq a^{k-1}\sum_{n<h\le k-1}\binom{k-1}{h}
\left(\frac ck\right)^h(q^{n-h}-1).\nonumber
\end{align}
We need an upper bound on $T_n/p_0$ and a lower bound on $H_n/p_0$.

For the first estimate, write $h=n+j$. Since $a^{k-1}/p_0<3$, $\binom{k-1}{n+j}k^{-n-j}\le1/(n+j)!$, and
\begin{align*}
q^{-j}-1\le\frac jn q^{1-j},\qquad
(n+j)!\ge(n+1)!(n+2)^{j-1},
\end{align*}
we obtain
\begin{align}\label{eq:appendix-tail-error}
\frac{T_n}{p_0}
&<\frac{3c^{n+1}}{n(n+1)!}
\sum_{j\ge1}j\left(\frac{c}{q(n+2)}\right)^{j-1}\nonumber\\
&=\frac{3c^{n+1}}{n(n+1)!}
\left(1-\frac{c}{q(n+2)}\right)^{-2}
<\begin{cases}
1/4000,&n=4,\\
1/(500n^3),&n\ge5.
\end{cases}
\end{align}
For $n=4$, replacing $c$ by $1/2$ gives $9/36980<1/4000$. For $n=5$, the corresponding bound is $49/3145728<1/(500\cdot5^3)$. For $n\ge5$, multiply the displayed bound by $n^3$; both factors
\begin{align*}
\frac{n^2}{2^{n+1}(n+1)!}
\quad\text{and}\quad
\left(1-\frac{n+1}{2n(n+2)}\right)^{-2}
\end{align*}
decrease with $n$, so the bound at $n=5$ proves the remaining cases.

For the lower bound on $H_n/p_0$, put
\begin{align*}
\tau\coloneqq\frac3{6n+4},\qquad
A_n\coloneqq1+\tau+\frac{\tau^2}{2}+\frac{\tau^3}{6},
\end{align*}
and define
\begin{align*}
B_{n,k}(c)&\coloneqq1+\frac{(k-1)\delta}{n}
+\frac{(k-1)(k-2)\delta^2}{2n^2},\\
C_k(c)&\coloneqq1+\frac{1/2-c}{k}+\frac1{12k^2}.
\end{align*}
These expressions come from truncating three positive series. First,
\begin{align*}
\log(1+u)\le\frac{u(6+u)}{6+4u}\qquad(u\ge0),
\end{align*}
with strict inequality for $u>0$: the derivative of the right-hand side minus the left-hand side is $u^3/((1+u)(2u+3)^2)$. Taking $u=1/n$ gives $1+n\log q>\tau$, and therefore $\mathrm e q^n>\exp(\tau)>A_n$.

Second, using $\log((1+u)/(1-u))>2u+2u^3/3$ for $0<u<1$ with $u=1/(2(k+c-1/2))$, we get
\begin{align*}
-\log p_0
&>\frac{k}{k+c-1/2}+\frac{k}{12(k+c-1/2)^3}\\
&>1+\frac{1/2-c}{k}+\frac1{12k^2}.
\end{align*}
Thus $\mathrm e^{-1}/p_0>C_k(c)$. Third, the binomial expansion gives $(1+\delta/n)^{k-1}\ge B_{n,k}(c)$. Multiplying these three inequalities yields
\begin{align}\label{eq:appendix-rational-lower-bound}
\frac{H_n}{p_0}>A_nB_{n,k}(c)C_k(c).
\end{align}

We compare this rational lower bound with the required value by defining
\begin{align*}
\Phi(n,k,c)\coloneqq A_nB_{n,k}(c)C_k(c)-1
-\widehat z_{n+1}\left(\frac{k}{k+c}\right).
\end{align*}
For $n\le k-2$, this expression increases with $c\in[11/24,1/2]$. To verify the assertion, differentiate $B_{n,k}(c)C_k(c)$. The derivative has the form $\alpha/n+\beta/n^2-1/k$, where
\begin{align*}
\alpha&=\frac{k-1}{(k+c)^2}
\left(k+\frac12-2c+\frac{1/12-c^2}{k}\right)>0,\\
\beta&=\frac{(k-1)(k-2)c}{(k+c)^3}
\left(k+\frac12-\frac{3c}{2}+\frac{1/12-c^2/2}{k}\right)>0.
\end{align*}
It is therefore enough to check its positivity at $n=k-2$. After multiplication by $2k(k-2)(k+c)^3$, the derivative there is
\begin{align}\label{eq:appendix-derivative-polynomial}
&(3-6c)k^3+\left(-\frac56+14c-15c^2\right)k^2\nonumber\\
&\qquad+\left(-\frac16-\frac{5c}{3}+21c^2-5c^3\right)k
+c\left(7c^2-\frac13\right)>0.
\end{align}
On the stated interval, the coefficient of $k^3$ is nonnegative, and the other three coefficients are at least $103/48$, $3325/1152$, and $7205/13824$, respectively. This proves the strict sign in \cref{eq:appendix-derivative-polynomial}. Finally, $\widehat z_t(u)$ increases with $u\in(0,1]$, since its derivative has positive numerator $19t^2+8tu-4u^2$. As $k/(k+c)$ decreases with $c$, the final term of $\Phi$ also increases with $c$.

The following two exact polynomial expansions complete the rational comparison. If $n\ge5$, put
\begin{align*}
d\coloneqq6n+4,\qquad K\coloneqq24k+11,\qquad U\coloneqq(n+1)K-24k,
\end{align*}
and define
\begin{align}\label{eq:appendix-big-polynomial}
P(n,k)\coloneqq
\frac5{124416}d^3k^2n^3K^2(n+1)U
\left(\Phi(n,k,11/24)-\frac1{500n^3}\right).
\end{align}
The prefactor is positive and clears all variable denominators, leaving a polynomial over $\mathbb Q$. To check its sign throughout $n\ge5$ and $k\ge n+2$, substitute $n=5+v$ and $k=7+v+r$, where $v,r\ge0$. Writing
\begin{align*}
P(5+v,7+v+r)=\sum_{0\le j\le5}r^jP_j(v),
\end{align*}
the exact coefficient calculation gives
\begin{center}
\begin{tabular}{c|rrrrrr}
$j$&0&1&2&3&4&5\\\hline
$\deg P_j$&12&11&10&9&8&7\\
smallest coefficient of $P_j$&6&25&40&30&10&1
\end{tabular}
\end{center}
Every coefficient is positive, including the constant coefficient of $P_0$. Hence $P(n,k)>0$ in the required range, and monotonicity in $c$ gives
\begin{align*}
\Phi(n,k,c_k)>\frac1{500n^3}\qquad(n\ge5,\ k\ge n+2).
\end{align*}

For $n=4$, use the stronger lower bound
\begin{align*}
c^\dagger\coloneqq\frac{11}{24}+\frac1{24k}<c_k,
\end{align*}
which follows from \cref{eq:appendix-crossing}. Put $K_4\coloneqq(3k+1)(8k+1)$ and $U_4\coloneqq5K_4-24k^2$. The expression
\begin{align}\label{eq:appendix-small-polynomial}
P_4(k)\coloneqq\frac{5488000}{56441883}k^2K_4^2U_4
\left(\Phi(4,k,c^\dagger)-\frac1{4000}\right)
\end{align}
is a polynomial of degree eight. The coefficients of $P_4(6+r)$, in increasing order of the power of $r$, are
\begin{align*}
\bigg(&\frac{181495655296977271}{10836841536},\,
\frac{686615422318476713}{43347366144},\,
\frac{90220597308501601}{14449122048},\,
\frac{28681346604690545}{21673683072},\\
&\frac{3542063783479817}{21673683072},\,
\frac{178426019505295}{14449122048},\,
\frac{28568345451247}{43347366144},\,
\frac{57181472051}{1806140256},\,1\bigg).
\end{align*}
They are all positive, so $P_4(k)>0$ for $k\ge6$ and
\begin{align*}
\Phi(4,k,c_k)>\frac1{4000}.
\end{align*}
The accompanying file \texttt{verification/verify\_certificates.py} expands \cref{eq:appendix-fourth-identity,eq:appendix-big-polynomial,eq:appendix-small-polynomial} over $\mathbb Q$, checks the displayed identities, and writes every coefficient of the last two polynomials to separate tables. These are exact rational calculations; the positivity tests do not use numerical sampling.

Combining these bounds on $\Phi$ with \cref{eq:appendix-tail-error,eq:appendix-rational-lower-bound}, we obtain
\begin{align*}
z_t=\frac{H_n}{p_0}-\frac{T_n}{p_0}-1>\widehat z_t(a).
\end{align*}
\Cref{eq:appendix-zhat,eq:appendix-target-reduction} now prove $r^*(W_t)-(19/20)M^*(W_t)>x/5$ for $5\le t<k$.

At the last node $t=k$, the definitions give $1-V_k=(1-x)^{k-1}$. Substitution into \cref{eq:appendix-node-identities} yields $r^*(W_k)=M^*(W_k)=1$. Hence
\begin{align*}
r^*(W_k)-\frac{19}{20}M^*(W_k)=\frac1{20}>\frac x5,
\end{align*}
since $k\ge6$. Between consecutive nodes, $V^*$ and $r^*$ are affine. We have already proved $a^k>V^*$ at the endpoints, so this inequality holds throughout each segment. The function
\begin{align*}
M^*=x+\frac{xp_0}{a^k-V^*}
\end{align*}
is convex on the segment. Therefore $r^*-(19/20)M^*$ is concave and stays strictly above $x/5$ between its endpoint values. Together with \cref{eq:appendix-rb-ratio} and the positivity of the interpolation coefficients, this proves every assertion of \cref{lem:comparison-estimates}.
